\chapter*{Vorwort}
\addcontentsline{toc}{chapter}{Vorwort}

\begin{merksatz}
\boxtitel{Wo die bisherigen Bände aufhören}
Band 19 fragte, wie man das Beste findet. Die Antwort setzte etwas voraus, das dort unausgesprochen blieb: dass die Welt nicht zurückrechnet.

\begin{center}
{\small
\renewcommand{\arraystretch}{1.5}
\begin{tabular}{@{}P{2.4cm}P{6.2cm}@{}}
\toprule
\textbf{Optimierung} & Man wählt unter Nebenbedingungen, die feststehen.\\
\textbf{Spieltheorie} & Die Nebenbedingungen sind Entscheidungen anderer --- die ihrerseits die eigenen voraussehen.\\
\bottomrule
\end{tabular}}
\end{center}

\begin{center}
\begin{minipage}{0.9\textwidth}
\textbf{Sobald jemand anderes mitrechnet, gibt es kein Optimum mehr, sondern nur noch ein Gleichgewicht.}
\end{minipage}
\end{center}

\begin{center}
\begin{minipage}{0.88\textwidth}
\kursiv{Das ist ein Begriffswechsel, kein Verfahrenswechsel --- und er verändert alles. Ein Optimum ist die beste Wahl. Ein Gleichgewicht ist eine Lage, aus der niemand allein ausbrechen will. \textbf{Die beiden fallen im Allgemeinen nicht zusammen}, und der berühmteste Fall dieses Bandes ist einer, in dem das Gleichgewicht für alle Beteiligten schlechter ist als eine andere erreichbare Lage.}
\end{minipage}
\end{center}
\end{merksatz}

\begin{merksatz}
\boxtitel{Drei Ergebnisse, die man nicht erwartet}
\begin{center}
{\small
\renewcommand{\arraystretch}{1.55}
\begin{tabular}{@{}P{1.0cm}P{7.4cm}@{}}
\toprule
$1$ & In jedem endlichen Zweipersonen-Nullsummenspiel gibt es einen Wert --- eine Zahl, die beide Seiten sich gegenseitig garantieren können. Vorher weiß niemand, wie der andere zieht; nachher bereut keiner. \kursiv{(Kapitel 2)}\\
$2$ & Jedes endliche Spiel hat ein Gleichgewicht --- beliebig viele Spieler, beliebige Auszahlungen. Der Beweis benutzt keine Rechnung, sondern einen Fixpunktsatz. \kursiv{(Kapitel 6)}\\
$3$ & Es gibt kein Abstimmungsverfahren für drei oder mehr Möglichkeiten, das drei naheliegende Forderungen erfüllt. Nicht \enquote{keines ist bekannt} --- es gibt keines. \kursiv{(Kapitel 12)}\\
\bottomrule
\end{tabular}}
\end{center}

\begin{center}
\begin{minipage}{0.88\textwidth}
\kursiv{Das erste Ergebnis ist der Ausgangspunkt des Fachs, das zweite seine Grundlage, das dritte seine schärfste Grenze. \textbf{Zwischen dem ersten und dem dritten liegen sechsundzwanzig Jahre} --- und der Weg von der Garantie zum Unmöglichkeitsbeweis ist die Erzählung dieses Bandes.}
\end{minipage}
\end{center}
\end{merksatz}

\begin{merksatz}
\boxtitel{Die sechs Teile}
\begin{center}
{\small
\renewcommand{\arraystretch}{1.5}
\begin{tabular}{@{}P{1.4cm}P{2.8cm}P{3.8cm}@{}}
\toprule
\textbf{Teil} & \textbf{Thema} & \textbf{Kapitel}\\
\midrule
I & Entgegengesetzte Interessen & $1$--$3$\\
II & Gleichgewichte & $4$--$6$\\
III & Über die Zeit & $7$--$8$\\
IV & Evolution & $9$--$10$\\
V & Regeln entwerfen & $11$--$14$\\
VI & Stand & $15$\\
\bottomrule
\end{tabular}}
\end{center}

\schritt{Der Leitfall}
Wie in den Bänden 18 bis 20 bleibt Neustadt der Schauplatz --- diesmal mit einem zweiten Haus.

\begin{center}
{\small
\renewcommand{\arraystretch}{1.5}
\begin{tabular}{@{}P{1.4cm}P{6.8cm}@{}}
\toprule
I & Klinikum und Marienhospital wählen Notarztstandorte.\\
II & Beide entscheiden über den Ausbau einer Spezialabteilung --- und über den Einsatz von Reserveantibiotika.\\
III & Verhandlungen mit den Kassen, Jahr für Jahr.\\
IV & Resistenzen: Bakterien spielen ebenfalls, ohne es zu wissen.\\
V & Dienstplan, Organverteilung, Ausschreibung.\\
\bottomrule
\end{tabular}}
\end{center}

\begin{center}
\begin{minipage}{0.88\textwidth}
\kursiv{Teil~IV setzt unmittelbar bei Band 18, Kapitel~14 an. Dort hieß es, das Wettrüsten zwischen Therapie und Resistenz sei im Kern ein Spiel --- und der Beweis blieb offen. \textbf{Er wird in Kapitel~10 nachgeholt}, mit derselben Gompertz-Kurve und denselben Mutationsraten.}
\end{minipage}
\end{center}
\end{merksatz}

\begin{warnung}
\boxtitelr{Drei Warnungen}
\begin{center}
{\small
\renewcommand{\arraystretch}{1.55}
\begin{tabular}{@{}P{2.2cm}P{6.4cm}@{}}
\toprule
\textbf{Erste} & Die Auszahlungen sind das Modell. Wer sie falsch ansetzt, bekommt eine tadellose Lösung für ein anderes Spiel. \kursiv{Das ist die häufigste Fehlanwendung des ganzen Fachs.}\\
\textbf{Zweite} & Die Theorie beschreibt, was rationale Spieler tun würden. Ob Menschen es tun, ist eine empirische Frage --- und die Antwort lautet oft nein. \kursiv{Kapitel 15 behandelt, was daraus folgt.}\\
\textbf{Dritte} & Ein Gleichgewicht ist keine Empfehlung. Es sagt, wo man landet, nicht wo man landen sollte.\\
\bottomrule
\end{tabular}}
\end{center}

\begin{center}
\begin{minipage}{0.88\textwidth}
\kursiv{Die dritte Warnung ist die wichtigste und wird am häufigsten überhört. Das Gefangenendilemma aus Kapitel~5 hat genau ein Gleichgewicht, und es ist für beide Spieler schlechter als die Zusammenarbeit. \textbf{Wer daraus schließt, man solle nicht zusammenarbeiten, hat den Satz gelesen und nicht verstanden.}}
\end{minipage}
\end{center}
\end{warnung}

\begin{merksatz}
\boxtitel{Was dieser Band voraussetzt}
\begin{center}
{\small
\renewcommand{\arraystretch}{1.48}
\begin{tabular}{@{}P{3.2cm}P{1.6cm}P{2.8cm}@{}}
\toprule
\textbf{Werkzeug} & \textbf{Band} & \textbf{gebraucht in}\\
\midrule
Wahrschein\-lichkeit & $12$ & Kapitel 3, 9\\
Lineare Optimierung & $19$ & Kapitel 3\\
Konvexität, Dualität & $19$ & Kapitel 3, 6\\
Differential\-gleichungen & $18$ & Kapitel 10\\
Fixpunktsätze & \textcolor{akzent}{---} & Kapitel 6\\
\bottomrule
\end{tabular}}
\end{center}

\begin{center}
\begin{minipage}{0.88\textwidth}
\kursiv{Die letzte Zeile ist die einzige echte Neuerung: Der Satz von Brouwer wird in Kapitel~6 eingeführt, nicht bewiesen --- und das ist die erste von neun Auslassungen, die der Anhang tabelliert. \textbf{Alles andere wird entweder hier entwickelt oder in einem früheren Band nachgeschlagen}, und wo etwas vorausgesetzt wird, steht die Fundstelle daneben.}
\end{minipage}
\end{center}
\end{merksatz}
